\documentclass[letterpaper]{article}
\usepackage{geometry}

%% font stuff

\usepackage{hyperref}
\usepackage{xspace}

\urlstyle{sf}
\usepackage[T1]{fontenc}
\usepackage{libertine}
\usepackage[libertine]{newtxmath}
\usepackage[scaled=0.8]{beramono}

\usepackage{microtype}
\usepackage{cmap}

\usepackage{amsmath}

%% colors, styles

\usepackage{xcolor}

\newlength{\myline}
\setlength{\myline}{0.7pt} % pgf "thick"

\colorlet{bg}{white}
\colorlet{fg}{black}
\colorlet{shaded}{black!10}
\colorlet{shadeborder}{black!30}
\definecolor{hlbg}{HTML}{F5F5A4} 
\colorlet{hlfg}{black}
\colorlet{commentcolor}{black!60!white}
\definecolor{errorcolor}{HTML}{a91616}
\definecolor{linkcolor}{HTML}{4b804c}
\definecolor{bitcolor}{HTML}{a91616}

\pagecolor{bg}
\color{fg}

%% highlighting

\usepackage[auto,outline]{contour}

\contourlength{1.2\myline}
\newcommand{\hl}[1]{%
    \relax\ifmmode%
        {}%
        \contour{hlbg}{\textcolor{hlfg}{${} #1 {}$}}%
        {}%
    \else%
        \contour{hlbg}{\textcolor{hlfg}{#1}}%
    \fi%
}

%% codebox stuff 

\usepackage{varwidth}

\usepackage{pipetex}
\pipetexcommand{./codebox2tex.pl}

\definecolor{boxbordercolor}{HTML}{000000}
\definecolor{boxbgcolor}{HTML}{FFFFFF}

\newcommand{\codebox}[1]{%
    \begin{varwidth}{\linewidth}%
        \upshape%   no slant in definition/theorem statement!
        \begin{tabbing}%
            ~~~\=\quad\=\quad\=\quad\=\kill
            #1
        \end{tabbing}%
    \end{varwidth}%
}

\newcommand{\fcodebox}[1]{%
    \fboxsep=3pt%
    \fcolorbox{boxbordercolor}{boxbgcolor}{\codebox{#1}}%
}

\newcommand{\titlecodebox}[2]{%
    \fboxsep=0pt%
    \fcolorbox{boxbordercolor}{boxbordercolor!10!boxbgcolor}{%
        \begin{varwidth}{\linewidth}%
            \centering%
            \fboxsep=3pt%
            \colorbox{boxbordercolor!10!boxbgcolor}{#1} \\
            \colorbox{boxbgcolor}{\codebox{#2}}%
        \end{varwidth}%
    }%
}

\newcommand{\hltitlecodebox}[2]{%
    \fboxsep=3pt%
    \colorbox{hlbg}{%
        \titlecodebox{#1}{#2}%
    }%
}

\newcommand{\hlcodebox}[1]{%
    \fboxsep=3pt%
    \colorbox{hlbg}{%
        \fcodebox{#1}%
    }%
}

\newcommand{\procheader}[1]{\underline{#1}}
\newcommand{\mycomment}[1]{\textcolor{commentcolor}{\small\textsl{// #1}}}

%% library stuff and math stuff

\renewcommand{\L}{\mathcal{L}}
\newcommand{\lib}[1]{\mathcal{L}_{\textsf{\textup{#1}}}}
\newcommand{\outputs}{\Rightarrow}
\newcommand{\link}{\diamond}

\newcommand{\indist}{\approxeq}
\renewcommand{\gets}{\twoheadleftarrow}

\renewcommand{\le}{\leqslant}
\renewcommand{\leq}{\leqslant}
\renewcommand{\ge}{\geqslant}
\renewcommand{\geq}{\geqslant}

\newcommand{\pct}{\mathbin{\%}}
\renewcommand{\le}{\leqslant}
\renewcommand{\leq}{\leqslant}
\renewcommand{\ge}{\geqslant}
\renewcommand{\geq}{\geqslant}

\newcommand{\secpar}{\lambda}

%% bits

\newcommand{\bit}[1]{\ensuremath{\textcolor{bitcolor}{\texttt{\upshape #1}}}\xspace}
\newcommand{\bits}{\{\bit0,\bit1\}}

%% algorithms

\newcommand{\algorithm}[1]{\ensuremath{\textsf{\upshape#1}}\xspace}

\newcommand{\KeyGen}{\algorithm{KeyGen}}
\newcommand{\Enc}{\algorithm{Enc}}
\newcommand{\Encaps}{\algorithm{Encaps}}
\newcommand{\Sign}{\algorithm{Sign}}
\newcommand{\Dec}{\algorithm{Dec}}
\newcommand{\Decaps}{\algorithm{Decaps}}
\newcommand{\MAC}{\algorithm{MAC}}
\newcommand{\Verify}{\algorithm{Verify}}
\newcommand{\Share}{\algorithm{Share}}
\newcommand{\Reconstruct}{\algorithm{Reconstruct}}
\newcommand{\Filter}{\algorithm{Filter}}
\newcommand{\Birthday}{\algorithm{Birthday}}
\newcommand{\Start}{\algorithm{Start}}
\newcommand{\Respond}{\algorithm{Respond}}
\newcommand{\Finalize}{\algorithm{Finalize}}
\newcommand{\RSA}{\algorithm{RSA}}
\newcommand{\Commit}{\algorithm{Commit}}
\renewcommand{\Check}{\algorithm{Check}}
\newcommand{\Extract}{\algorithm{Extract}}
\newcommand{\SimTrans}{\algorithm{SimTranscript}}

\newcommand{\subname}[1]{\textsc{#1}}


\begin{document}

\title{
    The Joy of Typesetting Cryptography
    \\[2pt]
    \Large How to make it look like \emph{The Joy of Cryptography}
}
\author{Mike Rosulek, \url{https://joyofcryptography.com}}
\maketitle

This is a self-contained example of the formatting used in \emph{The Joy of Cryptography}.
Please consult the source code of this document.

\section*{A few convenient macros}

\begin{center}
\begin{tabular}{ll}
    \texttt{\textbackslash bit\{0\}} & $\bit{0}$ \\
    \texttt{\textbackslash bits} & $\bits$ \\
    \texttt{\textbackslash indist} & $\indist$ \\
    \texttt{\textbackslash gets} & $\gets$ \\
    \texttt{\textbackslash lib\{test\}} & $\lib{test}$ \\
    \texttt{\textbackslash hl\{highlighted\} text} & \hl{highlighted} text \\
    \texttt{math mode: \$X \textbackslash hl\{\textbackslash gets \textbackslash bits\textasciicircum\textbackslash secpar\}\$} & math mode: $X \hl{\gets \bits^\secpar}$ \\
\end{tabular}
\end{center}

\section*{``codebox'' syntax:}

The \verb+\begin{pipetex}+ environment passes its body to the \texttt{codebox2tex.pl} script, and renders the output as TeX.
    See \url{https://github.com/rosulek/pipetex.sty/}.

\begin{center}
    \fcolorbox{bitcolor}{bitcolor!10!white}{You need to compile with the \texttt{-shell-escape} option to call external scripts!}%
\end{center}

\noindent
Here is a brief overview of the \texttt{codebox2tex.pl} domain-specific language:

\begin{itemize}
    \item Inside \texttt{code} ... \texttt{end} you can write multiple lines vertically. 
        Outside of \texttt{code}, linebreaks don't matter.
        \begin{center}
            \begin{minipage}{0.3\linewidth}
            \begin{verbatim}
code
    K \gets \bits^\secpar
    Y := F(K,X)
    Z := F(K,Y)
end
A := 1
B := 2
            \end{verbatim}
            \end{minipage}
            \hspace*{0.5in}
        \begin{pipetex}
            code
                K \gets \bits^\secpar
                Y := F(K,X)
                Z := F(K,Y)
            end
            A := 1
            B := 2
        \end{pipetex}
        \end{center}

    \item \texttt{lib} ... \texttt{end} is similar, but draws a box.
        \texttt{lib} takes an optional argument.
        You can have more than one \texttt{lib}/\texttt{code} box (often separated by \texttt{\textbackslash link}, \texttt{\textbackslash indist}, etc):
        \begin{center}
            \begin{minipage}{0.3\linewidth}
            \begin{verbatim}
lib
    K \gets \bits^\secpar
    Y := F(K,X)
    Z := F(K,Y)
end
\link
lib \lib{test}
    R \gets \bits^\secpar
    U := R \oplus Z
end
\indist
\cdots
            \end{verbatim}
            \end{minipage}
            \hspace*{0.5in}
        \begin{pipetex}
            lib
                K \gets \bits^\secpar
                Y := F(K,X)
                Z := F(K,Y)
            end
            \link
            lib \lib{test}
                R \gets \bits^\secpar
                U := R \oplus Z
            end
            \indist
            \cdots
        \end{pipetex}
        \end{center}


    \item 
        Lines that contain no \$ are assumed to be entirely math (exception: lines like ``\texttt{else:}'').
        Lines that start with \texttt{//} are rendered as comments.
        \begin{center}
            \begin{minipage}{0.4\linewidth}
            \begin{verbatim}
lib
    // math:
    K \gets \bits^\secpar

    // mixture of math and non-math:
    add $K$ to the set $\mathcal{K}$
end
            \end{verbatim}
            \end{minipage}
            \hspace*{0.5in}
        \begin{pipetex}
            lib
                // math:
                K \gets \bits^\secpar

                // mixture of math and non-math:
                add $K$ to the set $\mathcal{K}$
            end
        \end{pipetex}
        \end{center}

    \item 
        Four spaces for indentation:
        \begin{center}
            \begin{minipage}{0.3\linewidth}
            \begin{verbatim}
lib
    K \gets \bits^\secpar

    if $K == \bit0^\secpar$:
        for $i=1$ to $n$:
            print $\bit{oops}$
end
            \end{verbatim}
            \end{minipage}
            \hspace*{0.5in}
        \begin{pipetex}
            lib
                K \gets \bits^\secpar

                if $K == \bit0^\secpar$:
                    for $i=1$ to $n$:
                        print $\bit{oops}$
            end
        \end{pipetex}
        \end{center}



    \item 
        Start a subroutine with \texttt{proc}.
        If the subroutine name contains only \texttt{[a-zA-Z\textbackslash .]} then it is rendered nicely (with \texttt{\textbackslash subname}), otherwise assumed to be math:
        \begin{center}
            \begin{minipage}{0.3\linewidth}
            \begin{verbatim}
lib
    proc otp.enc(M)
        K \gets \bits^n
        return $K \oplus M$

    proc \mathbb{H}^+(A)
        if $H[A]$ undefined:
            H[A] \gets \bits^n
        return $H[A]$
end
            \end{verbatim}
            \end{minipage}
            \hspace*{0.5in}
        \begin{pipetex}
            lib
                proc otp.enc(X)
                    K \gets \bits^n
                    return $K \oplus X$

                proc \mathbb{H}^+(A)
                    if $H[A]$ undefined:
                        H[A] \gets \bits^n
                    return $H[A]$
            end
        \end{pipetex}
        \end{center}
\end{itemize}




\section*{Sample (from homework)}

    Prove that the following encryption scheme has CPA security (real vs random) if $F$ is a secure PRF (PRP):
    \begin{center}
        \begin{pipetex}
            code
                \begin{aligned} \mathcal K &= \bits^\secpar \\ \mathcal M &= \bits^\secpar \\ \mathcal C &= \bits^{2\secpar} \end{aligned}
            end
            \qquad
            code
                proc \Enc(K,M)
                    R \gets \bits^\secpar
                    X := F(K,R)
                    Y := F(K,M)
                    Z := F(K,R \oplus Y)
                    return $X \| Z$
            end
        \end{pipetex}
    \end{center}
    In other words, prove that:
    \begin{center}
        \begin{pipetex}
            $\quad \text{if} \quad$

            lib
                K \gets \bits^\secpar

                proc prf.query(X)
                    return $F(K,X)$
            end
            \indist
            lib
                proc prf.query(X)
                    if $L[X]$ undefined:
                        L[X] \gets \bits^\secpar
                    return $L[X]$
            end

            $\quad \text{then} \quad$

            lib
                K \gets \bits^\secpar

                proc cpa.enc(M)
                    // $\Enc(K,M)$:
                    R \gets \bits^\secpar
                    X := F(K,R)
                    Y := F(K,M)
                    Z := F(K,R \oplus Y)
                    return $X \| Z$
            end
            \indist
            lib
                proc cpa.enc(M)
                    X \gets \bits^\secpar
                    Z \gets \bits^\secpar
                    return $X \| Z$
            end
        \end{pipetex}
    \end{center}

\end{document}
